% Abstract (DECISIONS.md §F1; at most 1920 characters). Every claim is stated in the body with its conventions.
% Rename of 2026-10-15 (research/paper-review/rename/RENAME.md §4, Paper 1 item 1): Solomonoff's operator-induction problem and
% solutions first (Solomonoff 2008, §4), then the versions of SOI; "SOI" alone is the one-tape version (setting.tex, after def:setting:sfi).
\begin{abstract}
Solomonoff posed operator induction, predicting the answer to a new question from an unordered set of question--answer pairs (supervised learning), and gave universal-prior solutions with a convergence theorem. We study several versions of Solomonoff operator induction (SOI) and how they compare and behave. In the one-tape version, our default, a universal machine reads the same fair coin flips on each input. Its law dominates that of every such machine and depends on the universal machine only up to a constant factor. On distinct inputs it agrees with the coin-free and Kolmogorov versions and Solomonoff's per-query solution up to the complexity of the input set, but there is no coding theorem for function programs. On every consistent sequence, SOI's cumulative log loss exceeds that of every computable random labeller by at most its complexity plus a constant; with i.i.d.\ inputs and a computable deterministic target $f$, the averaged predictor has risk $O(K(f)/n)$. The one-tape version cannot model fresh label noise on repeated inputs; the per-query solution can. On unnormalised predictive probabilities, SOI's cumulative loss is order-independent: on a random cyclic permutation of $\{0,1\}^d$ revealed backwards along its cycle it pays about $(d-\log_2e)2^d$ bits, and no order-independent predictor pays less in expectation, while Solomonoff induction on the interleaved sequence, with self-delimiting blocks, pays $O(1)$ bits of label loss in expectation. Hypotheses that see the list of previous inputs give the universal chronological semimeasure. A dominance theorem for computable reveal laws holds pathwise for deterministic laws and, under finiteness assumptions and up to the complexities of the horizon and of the input sets, in expectation for randomised ones. We also treat Sam Eisenstat's coin example, out-of-distribution behaviour, computable versions and neural-network Bayes.
\end{abstract}
